% Part: first-order-logic
% Chapter: completeness
% Section: introduction

\documentclass[../../../include/open-logic-section]{subfiles}

\begin{document}

\iftag{FOL}
      {\olfileid{fol}{com}{int}}
      {\olfileid{pl}{com}{int}}

\olsection{Pambuka}

Teorema kelengkapan iku salah siji asil kang paling dhasar ngenani
logika. Ana rong formulasi, lan kita bakal mbuktekake yen kalorone
ekuivalen. Ing formulasi kapisan, teorema iki nyatakake prakara dhasar
ngenani sesambungan antarane akibat semantis lan sistem
!!{derivation} kita: yen !!a{sentence}~$!A$ dadi akibat saka sawatara
!!{sentence}s~$\Gamma$, uga ana !!a{derivation} kang netepake
$\Gamma \Proves !A$. Mula, sistem !!{derivation} iku sakuwat-kuwate
tanpa mbuktekake prakara kang sejatine ora dadi akibat.

Ing formulasi kapindho, teorema iki bisa dirumusake minangka asil
anane model: saben himpunan !!{sentence}s kang konsisten bisa
dicukupi. Konsistensi iku gagasan ing teori bukti: tegese sistem
!!{derivation} kita ora bisa ngasilake !!{derivation}s tartamtu.
Nanging, apa mung amarga ora ana !!{derivation}s jinis tartamtu
saka~$\Gamma$ banjur mesthi ana
\iftag{FOL}{!!a{structure}~$\Struct{M}$}{!!{valuation}{~$\pAssign{v}$}}
kanthi $\iftag{FOL}{\Sat{M}}{\pSat{v}}{\Gamma}$? Cathetan koreksi sumber OLPL-739: syarat nyukupi himpunan premis ditrapake marang cabang struktur lan valuasi; ing sumber, kurung kurawal nempatake syarat iki mung ing cabang valuasi. Sadurunge
teorema kelengkapan pisanan dibuktekake---malah sadurunge ana sistem
!!{derivation} kang saiki digunakake---matematikawan Jerman gedhe,
David Hilbert, duwe panemu yen konsistensi teori matematika njamin
anane objek kang dirembug teori mau. Ing layang marang Gottlob Frege,
dheweke nyatakake mangkene:
\begin{quote}
  Yen aksioma-aksioma apa wae kang diwenehake, bebarengan karo kabeh
  akibate, ora padha sulaya, aksioma-aksioma iku bener lan samubarang
  kang ditegesi dening aksioma kasebut ana. Kanggo aku, iki paugeran
  kabeneran lan anane.
\end{quote}
Frege mbantah kanthi tegas. Formulasi kapindho teorema kelengkapan
nuduhake yen Hilbert bener paling ora miturut teges iki: yen
aksioma-aksioma konsisten, mula ana \emph{sawijining}
\iftag{FOL}{!!{structure}}{!!{valuation}} kang ndadekake kabeh
aksioma mau bener.

Ana sebab liyane kang ndadekake teorema kelengkapan---utawa, luwih
trep, buktine---wigati. Teorema iki duwe sawetara akibat wigati, lan
sawetara bakal kita rembug dhewe. Contone, amarga saben !!{derivation}
kang nuduhake $\Gamma \Proves !A$ iku winates lan mula mung bisa
nggunakake !!{sentence}s kanthi cacah winates saka~$\Gamma$, saka
teorema kelengkapan nderek manawa yen $!A$ dadi akibat saka~$\Gamma$,
ukara iku wis dadi akibat saka sawijining subhimpunan winates
saka~$\Gamma$. Iki diarani \emph{kekompakan}. Kanthi ekuivalen, yen
saben subhimpunan winates saka~$\Gamma$ konsisten, mula $\Gamma$
dhewe mesthi konsisten.

Sanadyan teorema kekompakan nderek saka teorema kelengkapan kanthi
dalan ora langsung lumantar !!{derivation}s, \emph{bukti} teorema
kelengkapan uga bisa digunakake kanggo mbuktekake kekompakan kanthi
langsung. Cathetan koreksi sumber OLPL-044: sumber Inggris ngulang
artikel the ing frasa the proof of; terjemahan iki mulihake frasa
tunggal kang dikarepake. Bukti kasebut njupuk himpunan !!{sentence}s
kanthi sipat tartamtu---konsistensi---lan mbangun !!a{structure} saka
himpunan iki kang nduweni sipat tartamtu (ing perkara iki, struktur
kasebut nyukupi himpunan mau). Konstruksi kang meh padha bisa
digunakake kanggo mbuktekake kekompakan kanthi langsung, kanthi miwiti
saka himpunan !!{sentence}s kang saben subhimpunan winatese bisa
dicukupi tinimbang saka himpunan kang konsisten.
\iftag{FOL}{Konstruksi kasebut uga ngasilake akibat liyane, umpamane
  saben himpunan !!{sentence}s kang bisa dicukupi duwe model winates
  utawa !!{denumerable}s. (Asil iki diarani teorema
  L\"owenheim--Skolem.) Umume, konstruksi !!{structure}s saka
  himpunan !!{sentence}s asring digunakake ing logika, lan sok-sok uga
  ing filsafat.}{}

\end{document}
